% Figure dem_accuracy: generated by figures/rse_fig_dem_accuracy.py; placed where the text first discusses it.
\begin{figure}[!tbp]
  \centering
  \includegraphics[width=\textwidth]{../figures/out/rse_fig_dem_accuracy.pdf}
  \caption{\textbf{Accuracy of the waterline elevation model against ICESat-2.}
  (a)~The waterline elevation model (height on the DAHITI reference, \metric{params.dem_res_m}~m, shown on a \metric{params.fig_map_res_m}~m display grid) on the emerged seabed, with the \metric{h2.atl08.n_test_blocks} test blocks of \metric{h2.block_size_km}~km that hold ICESat-2 points (vermilion squares). Checked sectors are named in bold.
  (b)~Model height against the ICESat-2 ATL08 height (orthometric height plus the datum constant of \metric{h2.atl08.datum_constant_m}~m fitted on training blocks) for the $n = \metric{h2.atl08.n_test_points}$ test-block points on bare emerged seabed; dashed line, one to one.
  Mean error \metric{h2.atl08.dem.bias_m}~m, normalised median absolute deviation (NMAD) \metric{h2.atl08.dem.nmad_m}~m and root-mean-square error \metric{h2.atl08.dem.rmse_m}~m.
  (c)~Error (model minus ICESat-2) per sector and over all test blocks.
  Black circles, the elevation model against ATL08 (filled, mean; open, median); grey squares, the same against ATL06; coloured markers, the mean error of each grid against ATL08, after the grid's vertical offset fitted on training blocks.
  Shaded, $\pm\metric{params.checked_tol_m}$~m: a sector counts as checked when all four statistics of the elevation model lie within it.
  A value beyond the axis is drawn on its edge and labelled.
  (d)~NMAD (\metric{params.nmad_scale} times the median absolute deviation from the median) of the elevation model (black) and of each grid against ATL08, on a logarithmic axis.
  Right, the number of ATL08 points and of test blocks per row; the west and interfluve sectors hold \metric{h2.atl08.n_blocks.west} test block each, so no interval can be estimated for them.}
  \label{fig:rse_dem_accuracy}
\end{figure}
